\documentclass{rvcepaper}
\newcommand{\orrow}{\multicolumn{6}{|c|}{\bfseries OR}\\ \hline}
\newcommand{\instrA}[1]{\begin{seeinstr}\item Answer all questions from Part A. Part A questions should be answered in first three pages of the answer book only.\item Answer FIVE full questions from Part B. In Part B question number 2 is compulsory. Answer any one full question from 3 and 4, 5 and 6, 7 and 8, 9 and 10.#1\end{seeinstr}}
\begin{document}

%%%%%%%%%%%%%%%%%%%%%%%%%%%%%%%%%%%%%%%% ENTREPRENEURSHIP AND INTELLECTUAL PROPERTY RIGHTS
% ---------------- HS361TA, July/August 2026
\begin{seepaper}
\paperinfo{code={HS361TA},examline={VI Semester B.E. Regular / Supplementary Examination July / August 2026},dept={},title={ENTREPRENEURSHIP AND INTELLECTUAL PROPERTY RIGHTS},time={03 Hours},maxmarks={100}}
\seehead
\instrA{}
\seepartA
\begin{seeA}
\qaq{1}{1.1}{State any two scopes of entrepreneurship.}{02}{2}{1}
\qa{1.2}{How does entrepreneurship contribute to national development?}{02}{1}{1}
\qa{1.3}{Differentiate between an entrepreneur and an intrapreneur.}{02}{1}{1}
\qa{1.4}{Outline the concept of break-even analysis.}{02}{2}{3}
\qa{1.5}{What factors should an entrepreneur consider before choosing a source of finance?}{02}{2}{2}
\qa{1.6}{Mention any two objectives of Trademark protection.}{02}{2}{4}
\qa{1.7}{Give an example of a non-registrable mark.}{02}{2}{4}
\qa{1.8}{What is the validity period of a registered Design in India?}{01}{1}{2}
\qa{1.9}{What happens if a Trade Secret becomes public?}{01}{2}{2}
\qa{1.10}{Mention any two forms of intellectual property rights available to businesses.}{02}{1}{3}
\qa{1.11}{In the context of ideation, what is the primary role of integrating engineering principles?}{02}{2}{2}
\end{seeA}
\seepart{PART-B}
\begin{seeB}
\qb{2}{a}{A textile business owner continued using traditional methods despite competitors adopting automated production systems. After observing the success of automation over several years, he gradually introduced automated machines.
\begin{rom}\item Identify the type of entrepreneur.\item State two characteristics of this entrepreneur.\item Analyze the advantages and disadvantages of such entrepreneurial behavior.\end{rom}}{08}{3}{2}
 & b & Ravi works as a project manager in a manufacturing company. He developed an AI-based production monitoring system that improved efficiency by 25\%. The company provided resources and support for implementation.
\begin{rom}\item Is Ravi an entrepreneur or an intrapreneur? Justify your answer.\item Differentiate between entrepreneur and intrapreneur with respect to this case.\item Evaluate the benefits gained by the company through Ravi's initiative.\end{rom} & 08 & 4 & 2\\ \hline
\qb{3}{a}{A startup wants to launch an electric bicycle rental service in Bengaluru. Describe the steps involved in conducting market feasibility and demand analysis.}{08}{4}{2}
 & b & Rationalize the process of developing a Business Model and a Minimum Viable Product (MVP). How do these help in validating a startup idea? & 08 & 3 & 3\\ \hline
\orrow
\qb{4}{a}{Apply Porter's Generic Competitive Strategies to a supply-chain business and justify which strategy would help it achieve a competitive advantage. Support your answer with suitable examples.}{08}{3}{4}
 & b & Analyze how cross-disciplinary collaboration enhances technological innovation in startups. & 08 & 2 & 2\\ \hline
\qb{5}{a}{Describe the major financial statements and explain their importance in business decision-making.}{08}{2}{3}
 & b & Discuss the significance of contracts and corporate governance in entrepreneurial ventures. & 08 & 2 & 3\\ \hline
\orrow
\qb{6}{a}{Develop a marketing mix (4Ps) for a newly launched technology-based startup product.}{08}{3}{1}
 & b & A startup observed that some employees consistently exceeded performance targets while others struggled to meet expectations. The company had no formal appraisal system. Discuss how performance evaluation can improve employee productivity and organizational effectiveness. & 08 & 4 & 3\\ \hline
\qb{7}{a}{Discuss the functions of Trade Marks and Eco-Labels.}{06}{3}{4}
 & b & Define ``Patent Transfer''. Discuss the various methods by which patent rights can be transferred. & 10 & 3 & 4\\ \hline
\orrow
\qb{8}{a}{An international company develops a method to extract active compounds from a plant traditionally used in Ayurveda, claiming higher efficiency and stability. Can this method be patented, and how does it relate to the protection of traditional knowledge?}{08}{3}{4}
 & b & A software company patents an algorithm for data encryption. Another company independently develops a similar algorithm and uses it in its products. Can the first company claim infringement? & 08 & 4 & 4\\ \hline
\qb{9}{a}{Explain the features of Industrial Design and discuss its importance in product development and marketing.}{08}{2}{2}
 & b & A furniture manufacturer creates an innovative chair design but does not register it under the Designs Act. A competing company later registers a similar design and claims exclusive ownership.
\begin{rom}\item Identify the Intellectual Property involved.\item Discuss the legal issues faced by the original manufacturer.\item Explain why design registration is important.\item What lessons can businesses learn from this case?\end{rom} & 08 & 4 & 3\\ \hline
\orrow
\qb{10}{a}{A college computer laboratory installs one licensed software package on fifty computers without purchasing additional licenses. Discuss the legal consequences.}{08}{4}{4}
 & b & Explain the legal framework governing Trade Secret Protection in India. Discuss the challenges in protecting Trade Secrets. & 08 & 2 & 2\\ \hline
\end{seeB}
\end{seepaper}

% ---------------- HS361TA, August 2025
\begin{seepaper}
\paperinfo{code={HS361TA},examline={VI Semester B.E. Regular Examinations August -- 2025},dept={},title={ENTREPRENEURSHIP AND INTELLECTUAL PROPERTY RIGHTS},time={03 Hours},maxmarks={100}}
\seehead
\instrA{}
\seepartA
\begin{seeA}
\qaq{1}{1.1}{What are the two advantages of being an Intrapreneur with a company?}{02}{2}{1}
\qa{1.2}{Define Imitative (Adoptive) Entrepreneur.}{02}{2}{1}
\qa{1.3}{Define Financial Feasibility Analysis.}{02}{2}{2}
\qa{1.4}{Define Prototype Development.}{02}{2}{2}
\qa{1.5}{What is the role of CRM in building strong customer relationships and enhancing satisfaction levels?}{02}{2}{2}
\qa{1.6}{How do startups benefit from venture capital investment?}{02}{2}{3}
\qa{1.7}{Define Compulsory Licensing.}{02}{2}{4}
\qa{1.8}{Define Eco-label.}{02}{2}{4}
\qa{1.9}{How do performer's rights differ from copyright?}{02}{3}{4}
\qa{1.10}{Differentiate between industrial design and utility patent with example.}{02}{2}{5}
\end{seeA}
\seepart{PART-B}
\begin{seeB}
\qb{2}{a}{Anjali, a 45-year-old teacher, wants to start her own online learning platform. Her friends discourage her, saying she's too old, doesn't have a tech background, and won't be able to compete with big companies. They also tell her she needs a huge investment to start. Questions:
\begin{rom}\item Identify and describe any three myths reflected in the friends' opinions.\item What advice would you give Anjali to overcome these misconceptions?\item Mention two real-life entrepreneurs who succeeded despite similar challenges.\end{rom}}{08}{3}{1}
 & b & How can engineering colleges promote entrepreneurship among students? Propose initiatives or program that link technical skills with business development. & 08 & 2 & 1\\ \hline
\qb{3}{a}{Discuss the significance of break-even analysis for startups and emerging businesses.}{08}{4}{2}
 & b & Develop a detailed business plan framework for a new e-commerce startup, incorporating key elements such as market research, financial projections, and operational strategies. & 08 & 4 & 2\\ \hline
\orrow
\qb{4}{a}{You are part of a student team developing a tech-based startup idea. Discuss how you would approach the process of transforming this idea into a viable business opportunity by addressing the following aspects:
\begin{rom}\item Identifying market opportunities and analyzing current trends.\item Applying engineering principles during the ideation and design phase.\item Promoting cross-disciplinary collaboration for innovation.\item Conducting market feasibility and demand analysis.\item Evaluating technical feasibility through prototype development and proof of concept.\item Analyzing financial feasibility, including cost estimation and revenue projections.\end{rom}}{10}{2}{2}
 & b & Define SWOC analysis and discuss its significance in strategic planning. How can organizations use SWOC analysis to formulate effective business strategies? & 06 & 2 & 2\\ \hline
\qb{5}{a}{A company wants to launch a new energy drink. Explain how it can use market segmentation, targeting, and positioning to ensure successful market entry.}{08}{4}{3}
 & b & You are developing a startup offering an innovative eco-friendly packaging solution. Explain how you would create a compelling Unique Value Proposition (UVP) that clearly communicates your product's value to customers. In your answer, discuss:
\begin{rom}\item What UVP is and why it is important\item The steps you would take to develop a UVP\item How to differentiate your product from competitors\item How your UVP would influence branding, marketing and customer perception.\end{rom}
Support your response with logical reasoning and, where possible with relevant examples. & 08 & 4 & 3\\ \hline
\orrow
\qb{6}{a}{``Entrepreneurs face several legal and ethical challenges while establishing and running their ventures. Discuss the importance of understanding and managing the following aspects in entrepreneurship:
\begin{rom}\item Intellectual Property Rights (IPR) -- including patents, trademarks, and copyrights\item Contracts -- their role, essential elements, and how they protect business interests\item Corporate Governance -- principles, ethical considerations, and its impact on business sustainability\end{rom}}{08}{2}{3}
 & b & Examine the concept of Corporate Social Responsibility (CSR) and discuss its significance for entrepreneurs in ensuring ethical business practices. Include real-world examples to illustrate the impact of CSR on business success and reputation. Importance of CSR for Entrepreneurs. & 08 & 3 & 3\\ \hline
\qb{7}{a}{Explain the criteria for patentability of an invention and describe the step-by-step procedure to file a patent application.}{08}{3}{4}
 & b & Discuss the relationship between traditional knowledge and patent law. Explain the challenges involved in protecting traditional knowledge through patents and describe alternative ways to safeguard such knowledge. & 08 & 3 & 4\\ \hline
\orrow
\qb{8}{a}{Explain the concept of passing off and trademark infringement. Discuss the legal remedies available to business affected by these issues. Support your answer with relevant case studies to illustrate how courts have dealt with passing off and trademark infringement disputes.}{08}{3}{4}
 & b & Describe the absolute grounds for which a trademark registration can be denied. & 08 & 3 & 4\\ \hline
\qb{9}{a}{Discuss the legal tools available in India for protecting trade secrets. How do these tools help business safeguard their confidential information?}{08}{3}{3}
 & b & Explain the concept of passing off and trademark infringement. Discuss the legal remedies available to businesses affected by these issues. Support your answer with relevant case studies to illustrate how courts have dealt with passing off and Trade Mark. & 08 & 3 & 5\\ \hline
\orrow
\qb{10}{a}{Examine how copyright law strikes a balance between protecting creators' rights and serving the public interest. What challenges arise in enforcing copyright in today's digital environment?}{08}{3}{5}
 & b & Define industrial design and highlight its main characteristics. Outline the process to secure design protection and explain the reasons why design rights may be revoked. & 08 & 3 & 5\\ \hline
\end{seeB}
\end{seepaper}

% ---------------- 18HSI51, Dec-20/Jan-21
\begin{seepaper}
\paperinfo{code={18HSI51},examline={V Semester B.E. Dec-20/Jan-21 Examinations},dept={DEPARTMENT OF INDUSTRIAL ENGINEERING AND MANAGEMENT},title={18HSI51 / 61 -- INTELLECTUAL PROPERTY RIGHTS AND ENTREPRENEURSHIP},time={03 Hours},maxmarks={100}}
\seehead
\begin{seeinstr}\item Answer all questions from Part A. Part A questions should be answered in first three pages of the answer book only.\item Answer FIVE full questions from Part B. In Part B question number 2, 7 and 8 are compulsory. Answer any one full question from 3 and 4 \& one full question from 5 and 6.\end{seeinstr}
\seepartA
\begin{seeA}
\qaq{1}{1.1}{The broadcasting organizations like TV, Radio have a term of copyright for \blank[2cm].}{01}{--}{--}
\qa{1.2}{Plan of a building can be protected by \blank[2cm].}{01}{--}{--}
\qa{1.3}{A scientist involved in space research discusses a scientific principle with his colleague. The colleague uses the principle and invents a working device based on the principle. A third colleague who was instrumental in suggesting the manner of putting in place the arrangement of the mechanical devices claims himself to be the true and first inventor and applies for the patent. Is he entitled to apply for the patent? In case not, who out of the three, can apply for a patent.}{02}{--}{--}
\qa{1.4}{Collective Mark is registered by \blank[2cm].}{01}{--}{--}
\qa{1.5}{A scientist in DRDO develops anti-radar chemical which, when coated on aircrafts, would make them invisible from the screen of the radar. Before his employment he was made to sign a contract which provided that all the products of his research would belong to the employer that is the DRDO. Who in the given case is authorized to apply for the patent?}{02}{--}{--}
\qa{1.6}{Symbol of Maharaja of Air India is \blank[2cm].}{01}{--}{--}
\qa{1.7}{Imagine a cricketer sets up his own company to sell his own range of clothes. What type of intellectual property can he use to show that the clothes are made by his company?}{01}{--}{--}
\qa{1.8}{Does Trademarks assure quality of the product.}{01}{--}{--}
\qa{1.9}{Information Technology Amendment Act came into force on \blank[2cm].}{01}{--}{--}
\qa{1.10}{Are GM (genetically modified) crops are patentable?}{01}{--}{--}
\qa{1.11}{Computer program is considered as one of the \blank[2cm] Intellectual Property.}{01}{--}{--}
\qa{1.12}{The process of moving from concrete observations and stories to more abstract themes that span across groups of people is \blank[2cm].}{01}{--}{--}
\qa{1.13}{\blank[2cm] firm is associated the most with Design Thinking.}{01}{--}{--}
\qa{1.14}{Design Thinking typically helps in \blank[2cm].}{01}{--}{--}
\qa{1.15}{Caitlin is creating a new product for university students. She takes a design-thinking approach. Her first step is addressing who she is creating the product for and conducts research on understanding this target market. What is this step in the design thinking process?}{01}{--}{--}
\qa{1.16}{In order to register Industrial Design, Design should be new and \blank[2cm].}{01}{--}{--}
\qa{1.17}{Whether the shape and configuration of spring can be registered as Industrial Design in India? Analyze.}{01}{--}{--}
\qa{1.18}{\blank[2cm] is an important source of idea generation due to their familiarity with the needs of market?}{01}{--}{--}
\end{seeA}
\seepart{PART-B}
\begin{seeB}
\qb{2}{a}{Describe the inception and growth of World Intellectual property Organization by mentioning its role \& importance for protecting the new invention.}{08}{--}{--}
 & b & In order to be patentable, an invention must be worthy of a patent grant. The patent worthiness of an invention is assessed by the government through patentability requirements. Explain the essential requirements of patentability. & 08 & -- & --\\ \hline
\qb{3}{a}{Compulsory Licensing allows governments to license third parties (that is, parties other than the patent holders) to produce and market a patented product or process without the consent of patent owners. Discuss in brief.}{08}{--}{--}
 & b & Explain the different modes of Protection of Trade Secrets in relations to the role of Employees in Organizations, highlighting their main obligations as an employee. & 08 & -- & --\\ \hline
\orrow
\qb{4}{a}{Identify the grounds for refusal of registration of a Trademark in India. What kind of Trademarks can be registered under the (Indian) Trade Marks Act, 1999?}{08}{--}{--}
 & b & How does ``Likelihood of Confusion'' factor into Infringement of Trade Mark. Explain with examples of any two case situations. & 08 & -- & --\\ \hline
\qb{5}{a}{``Passing off: a fall back for potentially registerable marks that were not registered or failed to renew a registration''. Explain this statement with two case situations.}{08}{--}{--}
 & b & Copy Right protects the expression, not the ideas. Explain this statement with suitable illustrations. & 08 & -- & --\\ \hline
\orrow
\qb{6}{a}{Define the term Industrial Design and explain the essential features of Industrial Design.}{08}{--}{--}
 & b & Define Cyber Crimes. Explain the commonly categorized types of Cyber Crimes. & 08 & -- & --\\ \hline
\qb{7}{a}{Discuss the barriers which cause communication breakdown in an organizational setting.}{08}{--}{--}
 & b & Identify your own entrepreneurship style based on your personality traits, strengths, and weaknesses. & 08 & -- & --\\ \hline
\qb{8}{a}{Explain the method of Customer Centric selling skill in detail.}{08}{--}{--}
 & b & Identify and explain the central ideas \& principles of design thinking. & 08 & -- & --\\ \hline
\end{seeB}
\end{seepaper}

\end{document}
